\begin{figure}[htbp]
\centering
\begin{tikzpicture}[
  font=\small,
  hdr/.style={font=\bfseries, text=white, fill=VIUOrange, rounded corners=3pt,
    inner sep=5pt, minimum width=6.4cm, align=center},
  body/.style={draw=VIUDark!50, very thick, rounded corners=5pt, fill=VIULight,
    text width=6.0cm, inner sep=8pt, align=left, anchor=north},
  res/.style={draw=VIUOrange!75, thick, rounded corners=4pt, fill=VIUOrange!10,
    text width=6.0cm, inner sep=7pt, align=left, anchor=north, font=\footnotesize},
  bridge/.style={-{Latex[length=8pt,width=7pt]}, line width=1.6pt, draw=VIUDeep},
]
\def\colL{0}
\def\colR{8.9}

% -------- Cabeceras --------
\node[hdr] (HL) at (\colL,0) {Teoría base · Teorema~\ref{prop:convergencia}};
\node[hdr] (HR) at (\colR,0) {Sistema implementado · empírico};

% -------- Hipótesis / configuración --------
\node[body] (BL) at ($(HL.south)+(0,-0.18)$) {%
  $\bullet$~~Factor espacial $\Psi \equiv 1$\\[3pt]
  $\bullet$~~Comunicación global, $R = \infty$\\[3pt]
  $\bullet$~~Tareas homogéneas, $N = Kn$\\[3pt]
  $\bullet$~~Sin transición de modo};
\node[body] (BR) at ($(HR.south)+(0,-0.18)$) {%
  $\bullet$~~Descuento espacial $\Psi(d) < 1$\\[3pt]
  $\bullet$~~Comunicación local, $R$ finito\\[3pt]
  $\bullet$~~Transición reclutamiento\,$\to$\,transporte\\[3pt]
  $\bullet$~~Obstáculos y navegación reactiva};

% -------- Resultado / garantía --------
\node[res] (RL) at ($(BL.south)+(0,-0.32)$) {%
  \textbf{Juego potencial de congestión.}
  Equilibrio Nash $x_k^{*}=n/N$ con estabilidad local probada para
  $\beta \in (\beta_{\min},\beta_{\max})$.};
\node[res] (RR) at ($(BR.south)+(0,-0.32)$) {%
  \textbf{Sistema híbrido perturbado.}
  Convergencia práctica a una vecindad del equilibrio
  (Prop.~\ref{prop:perturbacion}); se caracteriza empíricamente.};

% -------- Puente de perturbación --------
\draw[bridge] (BL.east) -- (BR.west)
  node[midway, above, font=\footnotesize\itshape, text=VIUDeep, align=center]
    {perturbación\\acotada}
  node[midway, below, font=\scriptsize, text=VIUGray, align=center]
    {$\Psi(d),\,R,$ modo,\\obstáculos};

\end{tikzpicture}
\caption[Teoría base frente a sistema implementado]%
{Relación entre el caso base analizado formalmente (Teorema~\ref{prop:convergencia})
y el sistema implementado evaluado por simulación. El descuento espacial $\Psi(d)$, la
comunicación local, la transición de modo y los obstáculos actúan como una perturbación
acotada del juego potencial base; la Proposición~\ref{prop:perturbacion} delimita su
efecto y la convergencia práctica a una vecindad del equilibrio se caracteriza
empíricamente.}
\viuownsource
\label{fig:teoria_implementado}
\end{figure}
